\newpage
\thispagestyle{empty}

\noindent
\begin{minipage}[t]{0.43\textwidth}
    \centering
    {\fontsize{9pt}{10pt}\selectfont ĐẠI HỌC QUỐC GIA TP.HCM}\\
    {\fontsize{9pt}{10pt}\selectfont \textbf{TRƯỜNG ĐẠI HỌC BÁCH KHOA}}\\
    \rule{5cm}{0.4pt}
\end{minipage}
\hfill
\begin{minipage}[t]{0.53\textwidth}
    \centering
    {\fontsize{9pt}{10pt}\selectfont  \textbf{CỘNG HÒA XÃ HỘI CHỦ NGHĨA VIỆT NAM}}\\
    {\fontsize{9pt}{10pt}\selectfont  Độc Lập -- Tự Do -- Hạnh Phúc}\\
    \rule{5cm}{0.4pt}
\end{minipage}

\vspace{1cm}

\begin{center}
    \Large \textbf{NHIỆM VỤ ĐỒ ÁN TỐT NGHIỆP THẠC SĨ}
\end{center}

\vspace{0.5cm}

\fontsize{13pt}{15pt}\selectfont

% Personal information
\noindent

\begin{tabular}{l l}
\textbf{Họ và tên học viên:} Lê Phúc Đức & \textbf{Mã số học viên:} 2370116\\
\textbf{Ngày, tháng, năm sinh:} 29/09/2000 & \textbf{Nơi sinh:} Thành phố Hồ Chí Minh\\
\textbf{Chuyên ngành:} Khoa học máy tính & \textbf{Mã ngành:} 8480101 \\
\end{tabular}

\vspace{0.7cm}

\begin{flushleft}

\noindent \textbf{I. \hspace{1mm} TÊN ĐỀ TÀI:}

\textit{Tên Tiếng Việt:}

NGHIÊN CỨU VÀ PHÁT TRIỂN HỆ THỐNG TỰ ĐỘNG PHÁT HIỆN HIỆN TƯỢNG TRÔI DẠT VÀ CẬP NHẬT MÔ HÌNH HỌC MÁY THÍCH ỨNG

\textit{Tên Tiếng Anh:}

NGHIÊN CỨU VÀ PHÁT TRIỂN HỆ THỐNG TỰ ĐỘNG PHÁT HIỆN HIỆN TƯỢNG TRÔI DẠT VÀ CẬP NHẬT MÔ HÌNH HỌC MÁY THÍCH ỨNG

\end{flushleft}

\vspace{2mm}

\begin{flushleft}

\noindent \textbf{II. \hspace{1mm} NHIỆM VỤ VÀ NỘI DUNG:}

\end{flushleft}

\vspace{1mm}

\begin{longtable}{@{}p{4cm}p{\dimexpr\linewidth-4cm-4\tabcolsep\relax}@{}}
\textbf{Nhiệm vụ đồ án tốt nghiệp} & \textbf{Nội dung} \\
\endfirsthead
\textbf{Nhiệm vụ đồ án tốt nghiệp} & \textbf{Nội dung} \\
\endhead
Khảo sát và phân tích & Khảo sát hiện tượng trôi dạt khái niệm, các thống kê kiểm định, khoảng cách phân phối MMD và phát hiện drift không giám sát. \\
Đánh giá hiệu suất & Tiến hành đánh giá toàn diện các bộ phát hiện drift trên benchmark gồm 14 tập dữ liệu với 30 lần chạy độc lập. \\
\pagebreak
Phân loại trôi dạt & Xây dựng hệ thống SE-CDT để nhận dạng loại drift dựa trên hình dạng tín hiệu MMD không cần nhãn. \\
Thích ứng mô hình & Đề xuất và triển khai các chiến lược cập nhật mô hình thích ứng tương ứng với từng loại drift được phát hiện. \\
Tích hợp thời gian thực & Thiết kế và kiểm thử nguyên mẫu hệ thống xử lý dòng thời gian thực trên nền tảng Apache Kafka. \\
\end{longtable}

\begin{flushleft}
\noindent \textbf{III. \hspace{1mm} NGÀY GIAO NHIỆM VỤ:} 25/08/2025
\end{flushleft}

\vspace{2mm}

\begin{flushleft}
\noindent \textbf{IV. \hspace{1mm} NGÀY HOÀN THÀNH NHIỆM VỤ:} 17/06/2026
\end{flushleft}

\vspace{2mm}

\begin{flushleft}
\noindent \textbf{V. \hspace{1mm} CÁN BỘ HƯỚNG DẪN:}
\begin{itemize}
    \item \textbf{CBHD 1:} PGS. TS. Thoại Nam
    \item \textbf{CBHD 2:} TS. Nguyễn Quang Hùng
\end{itemize}
\end{flushleft}

\vspace{1cm}

% Signature block per BIEU MAU 5: CAN BO HUONG DAN on the left, the
% place/date line and TRUONG KHOA on the right, both annotated
% "(Ky va ghi ro ho ten)".
\noindent
\begin{minipage}[t]{0.45\textwidth}
\centering
\mbox{}\\[\baselineskip]
\textbf{CÁN BỘ HƯỚNG DẪN}\\
\textit{(Ký và ghi rõ họ tên)}
\vspace{3cm} % Space for signature
\end{minipage}
\hfill
\begin{minipage}[t]{0.5\textwidth}
\centering
\textit{Tp. HCM, ngày 17 tháng 6 năm 2026}\\[\baselineskip]
\textbf{TRƯỞNG KHOA KHOA HỌC \\
VÀ KỸ THUẬT MÁY TÍNH}\\
\textit{(Ký và ghi rõ họ tên)}
\vspace{3cm} % Space for signature
\end{minipage}
